\subsection{Summary}
This prospective cohort study used multi-omics to characterize gut microbiome changes during allo-HSCT and examine which changes preceded moderate-to-severe acute GVHD. The authors combined 16S rRNA sequencing, shotgun metagenomics, targeted and untargeted fecal metabolomics, machine-learning feature selection, and bacterial-metabolite network analysis. Across the transplant period, patients showed reduced bacterial richness and diversity, broad depletion of commensal \taxClostridiales, loss of microbiome-encoded functions involved in short-chain fatty acid production, reduced fecal acetate, propionate, and butyrate, and expansion of several \taxStaphylococcusSpecies. Patients who subsequently developed moderate-to-severe acute GVHD had more pronounced depletion of specific \taxClostridiales taxa, greater expansion of \taxEnterococcusFaecium and \taxParabacteroides, greater reductions in acetate and malonate, more extensive loss of microbiome-encoded functions, and more disrupted relationships between commensal bacteria and their metabolic products. In contrast, the overall decline in $\alpha$-diversity was similar between future GVHD and non-GVHD groups, indicating that changes in particular taxa, functions, metabolites, and cross-omic interactions were more informative than the global diversity change alone.
\par

The full 16S cohort comprised 172 samples across 105 allo-HSCT recipients. The first stool sample was collected, on average, approximately 4-5 days before stem-cell infusion, and the second was collected around day~+14, near engraftment. All analyzed post-transplant samples were collected before GVHD developed. The principal paired 16S analyses used 67 recipients with matched pre-transplant and engraftment samples, while paired shotgun metagenomics and metabolomic analyses were performed in smaller subsets. For the main GVHD comparison, patients who later developed grade II--IV acute GVHD were compared with patients who developed grade I or no GVHD. The authors calculated within-patient changes between the two time points and used these dynamics to identify taxonomic, functional, and metabolomic features associated with future GVHD; they also examined associations with GVHD grade, lower-gastrointestinal involvement, skin-only GVHD, and bacterial-metabolite interaction networks.
\par

\subsection{Identified microbial taxa}
Most GVHD-associated findings below refer to the change from the pre-transplant sample collected shortly before stem-cell infusion to the sample collected around engraftment, rather than abundance at either time point alone. Both samples preceded clinical GVHD.

\begin{itemize}

\item \taxClostridiales: Members of the \taxClostridiales order were broadly depleted after allo-HSCT, including prominently discussed genera such as \taxBlautia, \taxFaecalibacterium, and \taxOscillibacter. The depletion was more pronounced in patients who subsequently developed grade II--IV acute GVHD.
  \begin{itemize}
  \item \taxRuminococcaceae, \taxClostridiumIV, and \taxFlavonifractor were among the 16S-derived taxa that declined more in patients who later developed GVHD.
  \item Shotgun sequencing similarly identified greater depletion of \taxAnaeromassilibacillusSp, \taxButyricicoccusSp, and an \taxLachnoanaerobaculumSpecies in patients who later developed GVHD.
  \end{itemize}

\item \taxEnterococcusFaecium: \taxEfaecium expanded to a greater extent between the pre-transplant and engraftment samples in patients who subsequently developed moderate-to-severe acute GVHD.

\item \taxParabacteroides: \taxParabacteroides increased after allo-HSCT in patients who later developed GVHD but decreased in the comparison group. Greater expansion was also associated with more severe GVHD.

\item \taxStaphylococcusAureus, \taxStaphylococcusEpidermidis, and \taxStaphylococcusSimulans: These species expanded after allo-HSCT across the cohort and were not presented as specifically associated with future GVHD. Post-transplant \taxStaphylococcus abundance was inversely associated with fecal short-chain fatty acid levels, and the lower short-chain fatty acid concentrations measured after transplant were less inhibitory to \taxSaureus and \taxSepidermidis growth in vitro than the concentrations measured before transplant.

\item \taxFaecalibacterium, \taxClostridiumXIVb, and butyrate: \taxFaecalibacterium abundance was positively associated with butyrate before transplantation in both outcome groups. This relationship was lost after transplantation in patients who did not develop moderate-to-severe GVHD and reversed to a negative association in patients who did. Within the post-transplant GVHD interaction network, \taxClostridiumXIVb was negatively associated with \taxFaecalibacterium and positively associated with butyrate. These were bacterial-metabolite network findings rather than evidence that \taxClostridiumXIVb itself expanded in patients who developed GVHD.

\item GVHD severity and anatomical site:
  \begin{itemize}
  \item Greater depletion of \taxButyricicoccusBare showed a borderline association with more severe GVHD.
  \item Selected \taxBacteroides ASVs were positively associated with GVHD severity, but shotgun sequencing did not confirm a specific \taxBacteroidesSpecies as positively associated with severity.
  \item The pattern of greater \taxClostridiales depletion and greater \taxEfaecium and \taxParabacteroides expansion was most evident among patients who developed lower-gastrointestinal GVHD.
  \item \taxVeillonellaceae and its corresponding order \taxSelenomonadales showed distinct dynamics in the limited analysis of patients who developed skin-only GVHD.
  \end{itemize}

\end{itemize}
